\cotefile{156-7}
\scanpage{108}
\mspage{107}
pour
\[
\mathrm{omb}(X)^{\unsure{o|0|\circ}}\ne\varnothing\qquad\forall X\in M,
\]
car si on avait $\mathrm{omb}(X)^o=\varnothing$, on aurait $X\mathrel{\unsure{|\circ|\,|\,|o|}}X$, ce qui est exclu par Mag $5''$ (\unsure{antiréflexivité} de $|\circ|$). De plus, elle implique aussi At L2 (via $(***)$) et At L3 (via \deleted{Mag} Cor. de Mag $4''$).

Et l'axiome At C (\unsure{spec|sup}) sera automatiquement valable (alors qu'il ne l'est peut-être pas, et qu'il est \unsure{boiteux|boîteux}, dans le cas où on ne suppose pas $(****)$).

Inversement, si on part d'un $\unsure{\gtrless|\lessgtr}$-magasin, satisfaisant At L1, L2, L3 (pour pouvoir y définir raisonnablement $|\circ|$), la structure $\gtrless$ n'est autre que celle $\gtrless'$ définie par $|\circ|$, ssi $M$ satisfait At C (spec).

Donc le domaine commun des deux notions est, \unsure{résumé} par l'équivalence de catégories :
\diagram{\begin{tikzcd}[ampersand replacement=\&,row sep=large]
\text{$\gtrless$-magasins satisfaisant At L1, L2, L3 et At C (spec)} \arrow[d, Leftrightarrow] \\
\text{$|\circ|$-magasins satisfaisant $(****)$}
\end{tikzcd}}

Donc les $\gtrless$-magasins\inserted{(compatibiliseurs)} et les $|\circ|$-magasins\inserted{(\unsure{disparateurs})} apparaissent comme deux \deleted{types de} généralisations d'une même espèce de structure. Celle-ci paraît trop restrictive à mon sens, par la seule présence de At L1.
